% !TEX program = lualatex
% Overleaf安定版：機械加工・組立・溶接部品 要求仕様書
% Compiler: LuaLaTeX

\documentclass[a4paper,10pt]{ltjsarticle}

\usepackage[top=18mm,bottom=18mm,left=16mm,right=16mm]{geometry}
\usepackage{xcolor}
\usepackage{array}
\usepackage{enumitem}
\usepackage{fancyhdr}
\usepackage{lastpage}

% ===== 色：同系色のみ =====
\definecolor{SpecInk}{HTML}{1F2F3D}
\definecolor{SpecHead}{HTML}{2D4B63}
\definecolor{SpecSoft}{HTML}{EEF3F6}
\definecolor{SpecMuted}{HTML}{5F7180}

% ===== 全体設定 =====
\color{SpecInk}
\setlength{\parindent}{0pt}
\setlength{\parskip}{3.5pt}
\setlist[itemize]{
  labelindent=0pt,
  leftmargin=2.0em,
  labelwidth=1.1em,
  labelsep=0.5em,
  itemsep=2pt,
  topsep=2pt,
  parsep=0pt,
  partopsep=0pt
}
\setlist[enumerate]{
  labelindent=0pt,
  leftmargin=2.2em,
  labelwidth=1.3em,
  labelsep=0.5em,
  itemsep=2pt,
  topsep=2pt,
  parsep=0pt,
  partopsep=0pt
}

% ===== 本文幅：章見出し・表・箇条書きの基準幅 =====
% ヘッダ・フッタ罫線と同じ幅にするため、本文幅そのものを使用。
\newlength{\SpecBodyWidth}
\setlength{\SpecBodyWidth}{\linewidth}

% ===== 入力欄 =====
% 図番・仕様書番号等は、半角アンダースコアを避け、ハイフン推奨。
\newcommand{\DocTitle}{機械加工・組立・溶接部品 要求仕様書}
\newcommand{\SpecNo}{260523-01}
\newcommand{\SpecDate}{\today}
\newcommand{\ProjectName}{記入}
\newcommand{\PartName}{記入}
\newcommand{\DrawingNo}{記入}
\newcommand{\Quantity}{記入}
\newcommand{\DueDate}{記入}

% ===== ヘッダー・フッター =====
\pagestyle{fancy}
\fancyhf{}
\lhead{\small \DocTitle}
\rhead{\small 仕様書番号：\SpecNo}
\lfoot{\small\textcolor{SpecMuted}{\DocTitle}}
\rfoot{\small\textcolor{SpecMuted}{Page \thepage\ / \pageref{LastPage}}}
\renewcommand{\headrulewidth}{0.6pt}
\renewcommand{\footrulewidth}{0.4pt}

% ===== 見出し：ヘッダ・フッタ罫線と同じ幅 =====
\newcommand{\PageTitle}[1]{%
  \vspace*{2mm}%
  \noindent\colorbox{SpecHead}{%
    \parbox{\dimexpr\SpecBodyWidth-2\fboxsep\relax}{%
      \large\bfseries\textcolor{white}{#1}%
    }%
  }%
  \par\vspace{5mm}%
}

\newcommand{\BlockTitle}[1]{%
  \vspace{4mm}%
  \noindent\colorbox{SpecSoft}{%
    \parbox{\dimexpr\SpecBodyWidth-2\fboxsep\relax}{%
      \normalsize\bfseries\textcolor{SpecHead}{#1}%
    }%
  }%
  \par\vspace{1.5mm}%
}

\newcommand{\Note}[1]{%
  \vspace{1.5mm}%
  \noindent\fbox{%
    \parbox{\dimexpr\SpecBodyWidth-2\fboxsep-2\fboxrule\relax}{%
      \footnotesize\textcolor{SpecMuted}{#1}%
    }%
  }%
  \par\vspace{1.5mm}%
}

% ===== 表：全表を本文幅に統一 =====
% tabular* で外形幅を \SpecBodyWidth に固定。
% 各表の列幅指定が多少異なっても、表全体の幅は同一になる。
\newenvironment{SpecTable}[1]{%
  \par\vspace{1mm}%
  \noindent\footnotesize%
  \renewcommand{\arraystretch}{1.10}%
  \setlength{\tabcolsep}{3.0pt}%
  \begin{tabular*}{\SpecBodyWidth}{@{\extracolsep{\fill}}#1@{}}%
}{%
  \end{tabular*}%
  \par\vspace{1.5mm}%
}

\newcommand{\Head}[1]{\textbf{\textcolor{SpecHead}{#1}}}
\newcommand{\Lab}[1]{\textbf{\textcolor{SpecHead}{#1}}}
\newcommand{\YN}{要 / 不要}

\begin{document}

% 表紙は別途作成。本テンプレートは本文ページ用。

% =========================================================
\PageTitle{1. 情報}

\BlockTitle{1.1 基本情報}
\begin{SpecTable}{|p{28mm}|p{43mm}|p{30mm}|p{43mm}|}
\hline
\Lab{対象品名} & \PartName & \Lab{適用装置 / 案件名} & \ProjectName \\
\hline
\Lab{仕様書番号} & \SpecNo & \Lab{作成日} & \SpecDate \\
\hline
\Lab{希望数量} & \Quantity & \Lab{希望納期} & \DueDate \\
\hline
\Lab{要求者 所属・名} & 記入 & \Lab{依頼先} & 記入 \\
\hline
\end{SpecTable}

\BlockTitle{1.2 調達対象品}
\begin{SpecTable}{|p{9mm}|p{52mm}|p{35mm}|p{16mm}|p{34mm}|}
\hline
\Head{No.} & \Head{品名} & \Head{図面番号} & \Head{数量} & \Head{備考} \\
\hline
1 &  &  &  &  \\
\hline
2 &  &  &  &  \\
\hline
3 &  &  &  &  \\
\hline
4 &  &  &  &  \\
\hline
\end{SpecTable}


\clearpage

% =========================================================
\PageTitle{2. 依頼範囲・提示資料}

\BlockTitle{2.1 依頼範囲}
\begin{SpecTable}{|p{36mm}|p{24mm}|p{86mm}|}
\hline
\Head{項目} & \Head{要否} & \Head{指定内容 / 備考} \\
\hline
材料手配 & \YN &  \\
\hline
機械加工 & \YN &  \\
\hline
溶接 & \YN &  \\
\hline
組立 & \YN &  \\
\hline
熱処理・表面処理 & \YN &  \\
\hline
検査・検査成績書 & \YN &  \\
\hline
\end{SpecTable}

\BlockTitle{2.2 その他、依頼事項}
\begin{itemize}
  \item その他依頼事項：記入。
  \item 対象外作業：記入。
  \item 別途手配品：記入。
  \item 支給品：記入。
\end{itemize}

{\footnotesize\textcolor{SpecMuted}{不一致、疑義、製作上の制約を発見した場合は、製作前に発注者へ確認すること。}}

\BlockTitle{2.3 提示資料}
\begin{SpecTable}{|p{34mm}|p{58mm}|p{54mm}|}
\hline
\Head{分類} & \Head{資料名} & \Head{備考} \\
\hline
製作図 &  &  \\
\hline
3D CAD &  &  \\
\hline
部品表 &  &  \\
\hline
検査基準書 &  & 必要時 \\
\hline
\end{SpecTable}

\clearpage

% =========================================================
\PageTitle{3. 材料仕様・加工仕様}

\BlockTitle{3.1 材料仕様}
\begin{SpecTable}{|p{28mm}|p{43mm}|p{30mm}|p{43mm}|}
\hline
\Lab{材質} & 図面指示による / 下記指定 & \Lab{材料規格} & JIS / ASTM / EN / その他 \\
\hline
\Lab{ミルシート} & \YN & \Lab{代替材} & 不可 / 可 / 要承認 \\
\hline
\Lab{支給材} & 有 / 無 & \Lab{材料ロット管理} & \YN \\
\hline
\end{SpecTable}
\Note{代替材を提案する場合は、機械的性質、耐食性、熱処理性、溶接性、寸法安定性への影響を明記し、発注者の承認を得ること。}

\BlockTitle{3.2 部品別材料指定}
\begin{SpecTable}{|p{9mm}|p{48mm}|p{30mm}|p{30mm}|p{29mm}|}
\hline
\Head{No.} & \Head{部品名} & \Head{材質} & \Head{規格} & \Head{備考} \\
\hline
1 &  &  &  &  \\
\hline
2 &  &  &  &  \\
\hline
3 &  &  &  &  \\
\hline
\end{SpecTable}

\BlockTitle{3.3 機械加工の基本要求}
\begin{itemize}
  \item 図面指示寸法、公差、幾何公差、面粗さを満足すること。
  \item 図面に指示のないバリ、カエリ、鋭利エッジは除去すること。
  \item ねじ部、はめあい部、シール面、摺動面に損傷がないこと。
  \item 加工基準面、測定基準面を変更する場合は事前確認すること。
  \item 加工歪みが懸念される場合は、加工順序または応力除去の要否を確認すること。
\end{itemize}

\BlockTitle{3.4 一般公差・標準仕上げ}
\begin{SpecTable}{|p{42mm}|p{104mm}|}
\hline
\Lab{図面指示なき寸法公差} & 例：JIS B 0405-m 相当 / 個別指定 \\
\hline
\Lab{図面指示なき幾何公差} & 例：JIS B 0419 相当 / 個別指定 \\
\hline
\Lab{指示なき角部} & 糸面取り / C0.2から0.5程度 / 個別指定 \\
\hline
\Lab{指示なき面粗さ} & 例：Ra 6.3 / 個別指定 \\
\hline
\end{SpecTable}

\clearpage

% =========================================================
\PageTitle{4. 溶接仕様・組立仕様}

\BlockTitle{4.1 溶接仕様}
\begin{SpecTable}{|p{28mm}|p{43mm}|p{30mm}|p{43mm}|}
\hline
\Lab{溶接方法} & TIG / MIG / MAG / アーク / その他 & \Lab{溶接材料} & 図面指示 / 下記指定 \\
\hline
\Lab{外観要求} & 割れ、ブローホール、有害なアンダーカットなきこと & \Lab{歪み管理} & \YN \\
\hline
\Lab{溶接後処理} & スパッタ除去 / ビード仕上げ / 酸洗い & \Lab{非破壊検査} & 不要 / PT / MT / UT / RT \\
\hline
\end{SpecTable}

\BlockTitle{4.2 溶接に関する要求事項}
\begin{itemize}
  \item 溶接寸法、脚長、溶接範囲は図面指示によること。
  \item 溶接歪みにより、機能寸法、取付寸法、平面度、直角度を損なわないこと。
  \item 溶接後に機械加工を行う箇所は、必要な取り代を確保すること。
  \item 欠陥補修を行う場合は、補修方法を明確にし、必要に応じて承認を得ること。
\end{itemize}

\BlockTitle{4.3 組立仕様}
\begin{SpecTable}{|p{28mm}|p{43mm}|p{30mm}|p{43mm}|}
\hline
\Lab{組立範囲} & 全組立 / 部分組立 / 仮組 & \Lab{締結部品} & 支給 / 製作者手配 \\
\hline
\Lab{接着剤・固定剤} & 使用可 / 使用不可 / 指定品あり & \Lab{分解納入} & 可 / 不可 \\
\hline
\Lab{可動部確認} & \YN & \Lab{締付トルク管理} & 要 / 不要 / 指定部のみ \\
\hline
\end{SpecTable}

\BlockTitle{4.4 組立に関する要求事項}
\begin{itemize}
  \item 部品相互の干渉、芯ずれ、傾き、ガタつきがないこと。
  \item 締結部は必要なトルク、緩み止め、座面状態を満足すること。
  \item 可動部は、引っ掛かり、異音、過大抵抗がないこと。
  \item 組立後に分解が困難となる構造については、事前に確認すること。
\end{itemize}

\clearpage

% =========================================================
\PageTitle{5. 処理仕様・外観仕上げ}

\BlockTitle{5.1 熱処理・表面処理}
\begin{SpecTable}{|p{28mm}|p{43mm}|p{30mm}|p{43mm}|}
\hline
\Lab{熱処理} & なし / 焼入れ / 焼戻し / 焼鈍 / 応力除去 & \Lab{処理証明書} & \YN \\
\hline
\Lab{表面処理} & なし / 黒染め / 無電解Ni / アルマイト / 塗装 / その他 & \Lab{マスキング指定} & あり / なし \\
\hline
\Lab{硬度指定} & あり / なし & \Lab{処理後寸法管理} & 要 / 不要 / 重要部のみ \\
\hline
\end{SpecTable}

\BlockTitle{5.2 処理時の注意事項}
\begin{itemize}
  \item 処理により寸法、公差、ねじ、はめあい、面粗さを損なわないこと。
  \item 処理膜厚または熱変形が重要寸法に影響する場合は、製作前に確認すること。
  \item マスキング箇所、処理不可箇所は図面または別紙で明示すること。
  \item 処理後の錆、変色、剥離、膨れ、密着不良がないこと。
\end{itemize}

\BlockTitle{5.3 外観・仕上げ要求}
\begin{SpecTable}{|p{34mm}|p{112mm}|}
\hline
\Head{項目} & \Head{判定基準} \\
\hline
バリ・カエリ & 使用上有害なものがないこと。 \\
\hline
傷・打痕 & 機能面、シール面、摺動面に有害な傷・打痕がないこと。 \\
\hline
錆・変色 & 有害な錆、腐食、処理不良がないこと。 \\
\hline
汚れ・異物 & 切粉、粉塵、溶接スパッタ、過剰な油分がないこと。 \\
\hline
エッジ & 機能上必要な箇所を除き、鋭利でないこと。 \\
\hline
\end{SpecTable}

\BlockTitle{5.4 使用環境・特記事項}
\begin{itemize}
  \item 使用環境：記入。
  \item 温度、湿度、腐食環境：記入。
  \item 清浄度要求：記入。
  \item その他の禁止事項、注意事項：記入。
\end{itemize}

\clearpage

% =========================================================
\PageTitle{6. 検査要求}

\BlockTitle{6.1 検査項目}
\begin{SpecTable}{|p{48mm}|p{22mm}|p{24mm}|p{52mm}|}
\hline
\Head{検査項目} & \Head{要否} & \Head{記録提出} & \Head{備考} \\
\hline
外観検査 & \YN & \YN &  \\
\hline
主要寸法検査 & \YN & \YN &  \\
\hline
幾何公差検査 & \YN & \YN &  \\
\hline
ねじ検査 & \YN & \YN &  \\
\hline
表面粗さ検査 & \YN & \YN &  \\
\hline
硬度検査 & \YN & \YN &  \\
\hline
溶接外観検査 & \YN & \YN &  \\
\hline
非破壊検査 & \YN & \YN & PT / MT / UT / RT \\
\hline
組立動作確認 & \YN & \YN &  \\
\hline
\end{SpecTable}

\BlockTitle{6.2 重点管理寸法・特性}
\begin{SpecTable}{|p{8mm}|p{22mm}|p{44mm}|p{20mm}|p{28mm}|p{24mm}|}
\hline
\Head{No.} & \Head{図面位置} & \Head{寸法 / 特性} & \Head{公差} & \Head{検査方法} & \Head{備考} \\
\hline
1 &  &  &  &  &  \\
\hline
2 &  &  &  &  &  \\
\hline
3 &  &  &  &  &  \\
\hline
4 &  &  &  &  &  \\
\hline
5 &  &  &  &  &  \\
\hline
\end{SpecTable}
\Note{検査成績書が必要な場合は、測定値、使用測定器、検査日、検査者を記録すること。}

\BlockTitle{6.3 検査・判定に関する補足}
\begin{itemize}
  \item 検査範囲：全数、抜取、重要寸法のみ全数などを明記すること。
  \item 測定器指定：三次元測定機、粗さ計、硬度計など、必要時に指定すること。
  \item 判定保留時は、発注者確認後に処置すること。
\end{itemize}

\clearpage

% =========================================================
\PageTitle{7. 提出物・納入条件}

\BlockTitle{7.1 提出書類・納入物}
\begin{SpecTable}{|p{48mm}|p{22mm}|p{34mm}|p{42mm}|}
\hline
\Head{提出物} & \Head{要否} & \Head{提出タイミング} & \Head{備考} \\
\hline
見積書 & \YN & 発注前 &  \\
\hline
工程表 & \YN & 製作前 &  \\
\hline
材料証明書 & \YN & 納入時 &  \\
\hline
熱処理証明書 & \YN & 納入時 &  \\
\hline
表面処理証明書 & \YN & 納入時 &  \\
\hline
検査成績書 & \YN & 納入時 &  \\
\hline
非破壊検査記録 & \YN & 納入時 &  \\
\hline
完成写真 & \YN & 出荷前 / 納入時 &  \\
\hline
\end{SpecTable}

\BlockTitle{7.2 梱包要求}
\begin{itemize}
  \item 輸送中に傷、打痕、錆、変形が生じないよう梱包すること。
  \item 精密面、シール面、摺動面、ねじ部は保護すること。
  \item 防錆が必要な場合は、防錆油、防錆紙、乾燥剤等を使用すること。
  \item 重量物は安全に荷扱いできる荷姿とすること。
  \item 混載時は部品識別を明確にすること。
\end{itemize}

\BlockTitle{7.3 識別・納入条件}
\begin{itemize}
  \item 部品識別：品名、図面番号、Rev.、数量を表示すること。
  \item 個別刻印：要否および刻印位置を図面または別紙で指定すること。
  \item ロット管理：必要時はロット番号と対象数量を記録すること。
  \item 納品書：品名、図面番号、Rev.、数量を記載すること。
  \item 納入場所：記入。
\end{itemize}

\BlockTitle{7.4 納入時確認事項}
\begin{itemize}
  \item 納入品の品名、図面番号、Rev.、数量が発注内容と一致していること。
  \item 必要な提出書類が同梱または事前提出されていること。
  \item 梱包破損、部品の変形、錆、打痕、識別不備がないこと。
  \item 分納、代替品、仕様逸脱がある場合は、事前承認が記録されていること。
  \item 受入時に判定できない事項は、発注者と製作者で確認方法を協議すること。
\end{itemize}

\clearpage

% =========================================================
\PageTitle{8. 品質保証・変更管理}

\BlockTitle{8.1 不適合・疑義発生時の報告対象}
\begin{itemize}
  \item 図面または仕様を満足しない寸法、形状、外観。
  \item 材料、処理、溶接、組立に関する不適合。
  \item 納期に影響する問題。
  \item 図面・仕様の解釈に疑義がある事項。
  \item 製作上、仕様変更が必要と考えられる事項。
  \item 補修、再製作、特別採用が必要な事項。
\end{itemize}
\Note{不適合品の補修、再製作、特別採用は、発注者の承認なく実施しないこと。}

\BlockTitle{8.2 事前承認を要する変更}
\begin{itemize}
  \item 材料の変更。
  \item 加工方法、溶接方法、溶接材料の変更。
  \item 熱処理または表面処理条件の変更。
  \item 外注先、再委託先の変更。
  \item 図面・仕様からの逸脱。
  \item 納入形態または補修方法の変更。
\end{itemize}

\BlockTitle{8.3 見積時確認事項}
\begin{SpecTable}{|p{48mm}|p{40mm}|p{58mm}|}
\hline
\Head{確認項目} & \Head{回答} & \Head{備考} \\
\hline
製作可否 & 可 / 不可 / 条件付き可 &  \\
\hline
希望納期対応 & 可 / 不可 &  \\
\hline
材料手配 & 可 / 不可 &  \\
\hline
表面処理対応 & 自社 / 外注 / 不可 &  \\
\hline
溶接対応 & 自社 / 外注 / 不可 &  \\
\hline
検査成績書提出 & 可 / 不可 &  \\
\hline
\end{SpecTable}

\BlockTitle{8.4 受入条件}
\begin{enumerate}
  \item 図面および本仕様書の要求を満足していること。
  \item 必要な検査および提出書類が揃っていること。
  \item 外観、寸法、機能に重大な不適合がないこと。
  \item 梱包、識別、数量、Rev. が発注内容と一致していること。
  \item 不明点、逸脱、変更について事前承認が得られていること。
\end{enumerate}

\clearpage

% =========================================================
\PageTitle{9. 案件固有の特記事項}

\BlockTitle{9.1 設計・使用条件}
\begin{itemize}
  \item 使用環境：記入。
  \item 取付相手、インターフェース：記入。
  \item 荷重条件、使用条件：記入。
  \item 温度条件、耐食性要求：記入。
  \item 清浄度要求、禁止事項：記入。
  \item その他：記入。
\end{itemize}

\BlockTitle{9.2 最低記入項目}
\begin{itemize}
  \item 基本情報：品名、図面番号、Rev.、数量、納期。
  \item 技術情報：図面、材質、表面処理、熱処理、一般公差。
  \item 製作範囲：加工、溶接、組立、材料手配の要否。
  \item 品質要求：検査項目、検査成績書の要否、重点管理寸法。
  \item 納入条件：梱包、識別、提出書類、納入時確認事項。
  \item 変更管理：代替材、仕様逸脱、補修時の承認要否。
\end{itemize}

\clearpage

% =========================================================
\PageTitle{10. 改訂履歴}

\BlockTitle{10.1 改訂履歴}
\begin{SpecTable}{|p{16mm}|p{28mm}|p{74mm}|p{28mm}|}
\hline
\Head{Rev.} & \Head{日付} & \Head{改訂内容} & \Head{作成 / 承認} \\
\hline
0 & 記入 & 初版 & 記入 \\
\hline
 &  &  &  \\
\hline
 &  &  &  \\
\hline
 &  &  &  \\
\hline
\end{SpecTable}

\end{document}
